% Part: first-order-logic
% Chapter: introduction
% Section: sentences

\documentclass[../../../include/open-logic-section]{subfiles}

\begin{document}

\olfileid[fr]{fol}{int}{snt}

\olsection{\usetoken{P}{sentence}}

Nous disposons maintenant d'une ébauche de définition de la
satisfaction, c'est-à-dire de « vrai dans », pour les !!{structure}s et
les !!{formula}s. Cette définition fait toutefois intervenir un
élément supplémentaire, une assignation de !!{variable}s, alors que
nous voulions définir les !!{sentence}s. Comment nous débarrasser des
assignations, et que sont les !!{sentence}s ?

Vous vous souvenez probablement, depuis votre première initiation à la
logique formelle, d'une discussion sur la relation entre
!!{variable}s et quantificateurs. Un quantificateur est toujours suivi
d'!!a{variable}; dans la partie de l'!!{sentence} sur laquelle porte ce
quantificateur, sa « portée », nous considérons que la !!{variable} est
« liée » par ce quantificateur. La définition des !!{formula}s
n'exigeait pas que chaque !!{variable} possède un quantificateur
correspondant; les !!{variable}s dépourvues d'un tel quantificateur
sont « libres » ou « non liées ». Nous appellerons !!{sentence}s toutes
les !!{formula}s qui ne contiennent aucune !!{variable} libre.

L'idée intuitive qui permet de reconnaître si une occurrence
d'!!a{variable} dans !!a{formula}~$!A$ est liée ou libre, et quel
quantificateur la lie, n'est guère difficile à saisir. Vous avez
peut-être appris une méthode de vérification qui fait intervenir le
comptage des parenthèses. Nous devons néanmoins exiger une définition
précise. Comme nous avons défini les !!{formula}s par induction, nous
pouvons également définir par induction les occurrences libres et
liées d'!!a{variable}~$x$ dans !!a{formula}~$!A$. Pour notre langage
simplifié, la définition pourrait prendre la forme suivante :
\begin{enumerate}
  \item Si $!A$ est atomique, toutes les occurrences de $x$ qui y
  figurent sont libres; l'occurrence de $x$ dans~$\Atom{\Obj P}{x}$ est
  donc libre.
  \item Si $!A$ est de la forme $\lnot !B$, une occurrence de~$x$ dans
  $\lnot !B$ est libre si et seulement si l'occurrence correspondante
  de~$x$ est libre dans~$!B$; autrement dit, les occurrences libres de
  variables dans~$!B$ sont exactement les occurrences correspondantes
  dans~$\lnot !B$.
  \item Si $!A$ est de la forme $(!B \land !C)$, une occurrence de~$x$
  dans $(!B \land !C)$ est libre si et seulement si l'occurrence
  correspondante de~$x$ est libre dans~$!B$ ou dans~$!C$.
  \item Si $!A$ est de la forme $\lexists[x][!B]$, aucune occurrence
  de~$x$ dans~$!A$ n'est libre. Si $!A$ est de la forme
  $\lexists[y][!B]$, où $y$ est !!a{variable} différente de~$x$, une
  occurrence de~$x$ dans $\lexists[y][!B]$ est libre si et seulement
  si l'occurrence correspondante de~$x$ est libre dans~$!B$.
\end{enumerate}

Une fois les occurrences libres et liées des variables définies avec
précision, il suffit de poser qu'!!a{sentence} est !!a{formula} sans
occurrence libre d'!!a{variable}.

\end{document}
